\documentclass[11pt]{article}
\usepackage[T1]{fontenc}
\usepackage[utf8]{inputenc}
\usepackage{lmodern}
\usepackage[a4paper,margin=25mm]{geometry}
\usepackage{amsmath,amssymb,amsthm,mathtools,microtype}
\usepackage{booktabs}
\usepackage{hyperref}
\hypersetup{colorlinks=true,linkcolor=blue,citecolor=blue,urlcolor=blue,
pdftitle={The six-dimensional cross-polytope has signaling dimension three},
pdfauthor={},pdfsubject={Computer-assisted exact proof}}
\newtheorem{theorem}{Theorem}[section]
\newtheorem{lemma}[theorem]{Lemma}
\newtheorem{proposition}[theorem]{Proposition}
\newtheorem{corollary}[theorem]{Corollary}
\newtheorem{remark}[theorem]{Remark}
\newcommand{\R}{\mathbb R}
\newcommand{\supp}{\operatorname{supp}}
\newcommand{\sdim}{\operatorname{sign.dim}}
\newcommand{\cube}{\{-1,1\}^6}
\title{The six-dimensional cross-polytope has\protect\\
signaling dimension three}
\author{Anonymous}
\date{11 October 2026}

\begin{document}
\maketitle

\begin{abstract}
The signaling dimension of a state space is the least number of classical
states needed to simulate all its input--output correlations exactly, with
shared randomness. Recent bounds leave the six-dimensional cross-polytope
with signaling dimension either three or four. We prove, by an exact
computer-assisted argument, that the value is three. A primitive measurement
is encoded by an affinely independent weighted subset of the six-dimensional
hypercube. Supports of size at most five admit a direct three-state simulation.
After normalizing one vertex, the remaining six- and seven-point cases become
binary matrices modulo row and column permutations. Independent orbit
accounting reduces them to 14,303 and 118,827 classes. Exact rational filtering
leaves 39 and 178 valid rooted measurement records. Eighty-seven have the direct
small-facet simulation; for each of the other 130, the evidence archive gives
a rational probability distribution satisfying the complete finite
simulability criterion. The 217 records comprise 49 full cube-symmetry classes,
of which 28 need the general certificates. A verifier using only rational
arithmetic reads the supplied class lists, detects omitted measurements,
and checks every certificate; graph regeneration is optional.
An independent orbit-size audit proves that the representatives cover all
normalized supports. Combining the resulting upper bound with the known lower
bound gives the exact value.
\end{abstract}

\noindent\textbf{Keywords:} signaling dimension; cross-polytope; generalized
probabilistic theories; exact computation; symmetry reduction; rational
certificate.

\section{Introduction}

The signaling dimension of a physical state space is its exact classical
simulation cost: the smallest classical alphabet that reproduces every
transition matrix obtainable from the state space when shared randomness is
free. The notion originates in the study of hypersignaling and classical
simulation of quantum systems
\cite{NoHypersignaling,FrenkelWeiner,Frenkel,DallArnoTosiniBuscemi}.
It depends on the whole correlation set rather than on one discrimination or
coding task.

Exact values outside classical and quantum state spaces are scarce. Kai and
Dall'Arno developed an algorithm for rational polytopes and computed the
cross-polytope through affine dimension five, obtaining the value three in
dimensions three, four and five \cite{KaiDallArno}. Their general-purpose
calculation stopped there. Dall'Arno, Tosini and Buscemi gave the underlying
extremal-measurement and linear-programming framework together with a public
implementation \cite{DallArnoTosiniBuscemi}; Doolittle and Chitambar developed
separate signaling-polytope tools for certifying channel simulation costs
\cite{DoolittleChitambar}. No dimension-six cross-polytope result was found in
the inspected versions of these implementations \cite{SigdimSoftware,SignalingSoftware}; this is not an exhaustive
absence claim. Bencze and Frenkel subsequently proved that the
$d$-dimensional cross-polytope $\diamond_d$ satisfies
\begin{equation}\label{eq:seed-bounds}
 \left\lceil\log_2(d+2)\right\rceil
 \leq \sdim(\diamond_d)
 \leq \left\lfloor\frac{d+3}{2}\right\rfloor
\end{equation}
and gave a finite criterion tailored to its primitive measurements
\cite{BenczeFrenkel}. At $d=6$, the two bounds are three and four.

\begin{theorem}\label{thm:main}
The signaling dimension of the six-dimensional cross-polytope is
\[
 \sdim(\diamond_6)=3.
\]
\end{theorem}

Together with the known values in dimensions three through five and the lower
bound in \eqref{eq:seed-bounds}, Theorem~\ref{thm:main} identifies the complete
three-state range:
\begin{corollary}\label{cor:boundary}
For positive integers $d$,
\[
 \sdim(\diamond_d)=3 \quad\Longleftrightarrow\quad 3\leq d\leq6.
\]
\end{corollary}
\begin{proof}
The values for $3\leq d\leq5$ are known, and Theorem~\ref{thm:main} supplies
$d=6$. The one- and two-dimensional cross-polytopes have signaling dimension
at most two \cite{KaiDallArno}. For $d\geq7$, the lower bound in
\eqref{eq:seed-bounds} is at least four.
\end{proof}

The upper bound in Theorem~\ref{thm:main} is computer-assisted. The reduction
to a finite set is proved below. Every surviving linear feasibility instance
has a saved rational solution, so the scientific certificate does not depend
on floating-point optimization. Two replay programs check different possible
failure modes. One reclassifies all supports and verifies every inequality
using exact fractions. The other avoids the graph-generation software and
proves by orbit sizes that the supplied representatives are unique and
complete.

\section{Measurements and the finite criterion}

Write
\[
 \diamond_d=\left\{v\in\R^d:\sum_{i=1}^d|v_i|\leq1\right\}.
\]
An extremal effect of $\diamond_d$ is indexed by a hypercube vertex
$x\in\{-1,1\}^d$ and has the form
\[
 e_x(v)=\frac{1+x\cdot v}{2}.
\]
Thus a nonzero weight function $f:\{-1,1\}^d\to[0,1]$ represents a
measurement precisely when
\begin{equation}\label{eq:measurement}
 \sum_{x:x_i=1} f(x)=\sum_{x:x_i=-1}f(x)=1
 \qquad(1\leq i\leq d).
\end{equation}
Equivalently, $\sum_x f(x)=2$ and $\sum_x f(x)x=0$. The measurement is
primitive precisely when $\supp f$ is affinely independent
\cite[Section~2]{BenczeFrenkel}. It therefore has at most $d+1$ support
points.

We use the following specialization of the finite simulability criterion in
\cite[Proposition~2.1]{BenczeFrenkel}. For a finite set $U$, let
$\binom{U}{\leq r}$ denote its subsets of size at most $r$.

\begin{proposition}[Bencze--Frenkel criterion]\label{prop:criterion}
Let $f:\{-1,1\}^d\to[0,1]$ satisfy \eqref{eq:measurement}, and set
$U=\supp f$. Its transition matrix is simulable by an $r$-vertex simplex if
and only if there is a probability distribution $F$ on
$\binom{U}{\leq r}$ such that
\begin{equation}\label{eq:criterion}
 \sum_{x\in S}f(x)
 \leq \sum_{X:X\cap S\ne\varnothing}F(X)
\end{equation}
whenever $S\subseteq U$ lies in a coordinate facet of the hypercube.
\end{proposition}

The same distribution $F$ must work for all $2d$ input vertices, not a separate
distribution for each input. The supply--demand argument underlying the
criterion then supplies an input-dependent coupling with this common marginal.
Shared randomness draws a set $X$; the sender communicates one of its at most
$r$ outcomes. Thus three states mean a three-symbol alphabet, not three bits.
The empty subset may be omitted: moving its mass to a singleton only increases
the right-hand sides of \eqref{eq:criterion}.

Taking $S$ to be an entire coordinate-facet intersection gives demand one.
Consequently every $X$ with $F(X)>0$ must meet both sides of every coordinate.
This necessary transversal condition can prune candidate variables, but does
not replace the inequalities for proper subsets of facets.

\begin{lemma}[Small-facet construction]\label{lem:small}
Let $f$ be any measurement of $\diamond_d$. If a coordinate facet contains at
most $r-1$ points of its support, it is simulable by an $r$-vertex simplex.
In particular every measurement with at most five outcomes is three-state
simulable.
\end{lemma}
\begin{proof}
Choose such a facet $H$ and put $L=H\cap U$. The construction in the proof of
\cite[Theorem~2.2]{BenczeFrenkel} assigns
\[
 F(L\cup\{x\})=f(x)\qquad(x\in U\setminus H)
\]
and zero elsewhere. The opposite facet has total weight one, so $F$ is a
probability distribution supported on sets of size at most $r$.
If $S\cap L\ne\varnothing$, its coverage probability is one, which bounds
its facet demand. Otherwise its coverage probability is exactly
$\sum_{x\in S}f(x)$. These two cases prove \eqref{eq:criterion}.
For at most five outcomes, either side of some coordinate contains at most
two support points. This is the existing Bencze--Frenkel construction,
restated here; neither affine independence nor a new structural theorem is
required.
\end{proof}

\section{Binary-matrix reduction}

It remains to consider primitive supports of size six or seven. The next
lemma turns their enumeration into a small bicoloured-graph problem.

\begin{lemma}\label{lem:binary}
Choose $x_0\in U$ and apply coordinate sign changes so that
$x_0=(-1,\ldots,-1)$. For the other points write
\[
 x_j=-\boldsymbol 1+2b_j,
 \qquad b_j\in\{0,1\}^6\setminus\{0\},
\]
and let $B$ be the matrix with columns $b_j$. Then $U$ is affinely
independent if and only if $B$ has full column rank. It supports a primitive
measurement if and only if the system
\begin{equation}\label{eq:Bweights}
 Bw=\boldsymbol 1
\end{equation}
is consistent, has a unique solution, and that solution satisfies
\[
 w_j>0\quad\text{for all }j,
 \qquad w_0:=2-\boldsymbol 1^{\mathsf T}w>0.
\]
The measurement weights are $(w_0,w_1,\ldots)$.
\end{lemma}
\begin{proof}
Represent the affine function $e_x$ by the column $(1,x)$. Subtracting the
$x_0$ column from each other column replaces it by $(0,2b_j)$, proving the
rank assertion. If $w_0$ is the weight at $x_0$, the constant part of the
measurement equation gives $w_0+\sum_jw_j=2$. Its six linear parts become
\[
 -\left(w_0+\sum_jw_j\right)\boldsymbol 1+2Bw=0,
\]
which is exactly \eqref{eq:Bweights}. Full column rank makes a consistent
solution unique. Strict positivity is precisely the condition that all
listed points belong to the support.
\end{proof}

Permuting coordinates permutes the rows of $B$, while relabeling
$U\setminus\{x_0\}$ permutes its columns. A matrix with distinct nonzero
columns is the incidence matrix of a bicoloured graph. No valid matrix has a
zero row, because \eqref{eq:Bweights} has right-hand side $\boldsymbol1$.
Consequently, bicoloured graph isomorphism gives a lossless prefilter for the
remaining supports. Choosing a different $x_0$ can duplicate an orbit, but
cannot omit one. Accordingly the 217 retained records below are rooted
records, not full signed-coordinate symmetry classes.

There is also a standard cooperative-game interpretation. The columns of
$B$ are incidence vectors of coalitions on six players. Full column rank and
a positive solution of $Bw=\boldsymbol1$ characterize a minimal balanced
collection. Our additional condition $\sum_jw_j<2$ ensures positive weight
at the normalized vertex. Enumeration of minimal balanced collections is
established prior work \cite{MinimalBalanced}; our calculation does not claim
a new enumeration algorithm. The external catalogue is available on request,
and no direct catalogue-to-catalogue comparison is claimed here.

The program \texttt{genbg} from nauty 2.9.3 \cite{McKayPiperno} generated
the prefilter classes. Its output was audited independently. For
$k\in\{5,6\}$, the number of unordered $k$-subsets of distinct nonzero
binary columns that meet every row is
\begin{equation}\label{eq:labelled-count}
 N_k=\sum_{j=0}^6(-1)^j\binom6j
              \binom{2^{6-j}-1}{k}.
\end{equation}
An independent C++ program applies every one of the $6!=720$ row
permutations to each representative and sorts its columns. It found distinct
canonical keys and summed the resulting orbit sizes to $N_k$ in both cases.
Thus the supplied class lists are complete without assuming that
\texttt{genbg} itself returned a complete list.

\begin{table}[ht]
\centering
\caption{Exact rooted support enumeration and certificates. Rejected
prefilter records are not valid measurements.}
\label{tab:counts}
\begin{tabular}{lrr}
\toprule
 & six-point support & seven-point support \\
\midrule
Binary columns $k$ & 5 & 6 \\
Labelled supports $N_k$ & 6,054,006 & 63,602,770 \\
Coordinate-permutation classes & 14,303 & 118,827 \\
Rejected prefilter records & 14,264 & 118,649 \\
Valid rooted primitive measurements & 39 & 178 \\
Small-facet certificates & 24 & 63 \\
Exact criterion certificates & 15 & 115 \\
Infeasible valid measurements & 0 & 0 \\
Full cube-symmetry classes & 12 & 37 \\
Difficult full classes & 5 & 23 \\
\bottomrule
\end{tabular}
\end{table}

\section{Exact certificate computation}

Each graph representative was decoded into $B$. Equation
\eqref{eq:Bweights}, its consistency, rank and positivity were first used as a
floating-point producer prefilter. This left the 217 candidate valid classes
in Table~\ref{tab:counts}. The independent verifier then reclassified every
one of the 133,130 representatives using rational arithmetic, detected no
omitted valid class and obtained the same 217 classes. Thus completeness does
not depend on the producer's numerical gate. For 87 retained classes, some
coordinate facet contains at most two support points, and the distribution in
Lemma~\ref{lem:small} supplies a three-state simulation.

For each of the other 130 classes, all nonempty subsets $X\subseteq U$ with
$|X|\leq3$ were used as distribution variables. The constraints comprise
$F(X)\geq0$, $\sum_XF(X)=1$, and \eqref{eq:criterion} for every nonempty
subset of each of the twelve coordinate-facet intersections. Duplicate
subsets were removed. Exact linear rational arithmetic found every instance
feasible. The evidence manifest stores the nonzero values of one rational
$F$ for each class.

For a concrete example, record 5605 in the five-column group uses the binary
support encodings $(0,16,35,15,61,62)$ in that order: bit $i$ is coordinate
$i+1$, and a zero bit denotes sign $-1$. Every outcome has weight $1/3$.
The distribution gives mass $1/3$ to each of
\[
 \{1,2,3\},\qquad\{0,3,4\},\qquad\{0,2,5\}.
\]
These are subset masks $14,25,37$. Direct rational evaluation checks all
84 nonempty facet-subset occurrences (31 distinct constraints). The geometry
is essential: this distribution is not asserted to work on arbitrary supports.
The same shared draw is used for all twelve input columns.

Re-rooting every valid support at each of its vertices and minimizing over
all coordinate permutations gives 12 and 37 full signed-coordinate classes.
Of these, 5 and 23 are difficult. The catalogue maps every rooted record to
its full class and transports a rational simulator to canonical vertex order.
This reduces the explanatory catalogue to 28 difficult classes, without
replacing the exhaustive rooted coverage argument.

The independent verifier does not invoke the feasibility solver. Its default
mode reads the supplied class files, decodes every record separately, and uses a
local rational row-reduction implementation to recompute
\eqref{eq:Bweights}. It fails if any valid class is absent from the manifest.
For every saved distribution it checks nonnegativity, total mass one, and
every inequality \eqref{eq:criterion} by integer and rational arithmetic.
This replay found the same 39 and 178 valid classes and verified all 217
certificates. Regeneration with nauty is a separate optional provenance check;
it is not a prerequisite for this supplied-input replay. Completeness of the
supplied files is established separately by the orbit-size audit.

\begin{proof}[Proof of Theorem~\ref{thm:main}]
It suffices to simulate every primitive measurement
\cite{DallArnoTosiniBuscemi,BenczeFrenkel}. Such a measurement has at most
seven outcomes. Lemma~\ref{lem:small} handles supports of size at most five.
For sizes six and seven, Lemma~\ref{lem:binary} and the orbit audit show that
the class lists in Table~\ref{tab:counts} contain a representative of every
possible support. Exact filtering retains every possible primitive
measurement, and the rational certificates verify
Proposition~\ref{prop:criterion} with $r=3$ for each retained class.
Therefore $\sdim(\diamond_6)\leq3$.

The lower bound in \eqref{eq:seed-bounds} gives
$\sdim(\diamond_6)\geq\lceil\log_2 8\rceil=3$. Hence equality holds.
\end{proof}

\section{Reproducibility and limitations}

The evidence archive contains the two graph-class files, the 217 rational
certificates, the producer, the exact verifier, the independent orbit audit,
and hashes for every scientific input. The package's \texttt{MANIFEST.json}
binds final file sizes and SHA-256 digests; the AI index and verification
receipt identify the replay commands and their scope. Historical receipts
remain distinct from checks of the revised publication package.

The result is exact but computer-assisted. Its trust base includes the stated
source theorem, ordinary compiler and interpreter correctness, and inspection
of short replay programs. The orbit audit removes nauty completeness from the
trust base, although nauty remains the source of the candidate representatives.
The rational replay removes Z3, NumPy and SymPy from the certificate trust
base. A future structural proof could explain uniformly why the 28 difficult
full classes (130 rooted records) are feasible and would be preferable to
a class-by-class certificate. Different internal implementations remain
producer-side checking, not unaffiliated reproduction or formal verification.

The generic strategy of enumerating extremal measurements and solving finite
classical-simulation linear programs is prior work
\cite{DallArnoTosiniBuscemi,KaiDallArno}. The contribution here is the exact
dimension-six value, its three-state boundary consequence, and an assurance
chain that separately audits orbit coverage and rational feasibility. The
certificate instances themselves are small; the main computational burden is
complete support coverage. No new general signaling-dimension algorithm is
claimed.

The theorem resolves only $d=6$. It does not improve the general bounds in
\eqref{eq:seed-bounds}, nor does the computation determine the next open
dimensions.

\subsection*{Data and code availability}
All source, class lists, exact certificates, verification receipts and replay
instructions are included in the accompanying evidence archive. No external
dataset is used.
The versioned archive identifier is
\url{https://doi.org/10.5281/zenodo.23296390}.
The source repository is
\url{https://github.com/ipitchford/cross-polytope-signaling-dimension-six}.

\subsection*{Ethics declaration}
The work uses no human participants, animals, personal data or sensitive
data. No ethics approval was required.

\subsection*{Author contributions}
The scholarly attribution is Anonymous. OpenAI Codex assisted with recorded
searches, computation, checking and drafting. Publication authority and account
ownership do not establish a human mathematical contribution; none is inferred.

\subsection*{Conflict of interest}
No competing interests were reported in the supplied materials. This is not
independent verification of an anonymous creator's personal circumstances.

\subsection*{Funding}
No external funding was reported in the supplied materials. No institutional
affiliation or funding award is asserted.

\subsection*{AI disclosure}
OpenAI Codex assisted with literature screening, program construction,
computer-assisted proof search, exact replay and manuscript drafting. All
mathematical claims are intended to be assessed from the cited sources and
the archived certificates rather than from model authority.

\begin{thebibliography}{99}
\sloppy

\bibitem{MinimalBalanced}
D. Laplace Mermoud, M. Grabisch and P. Sudh\"olter,
``Minimal balanced collections and their application to core stability and
other topics of game theory,'' \emph{Discrete Applied Mathematics}
\textbf{341} (2023), 60--81.
\url{https://doi.org/10.1016/j.dam.2023.07.025}.
Accepted text: arXiv:2507.05898v1.

\bibitem{SigdimSoftware}
M. Dall'Arno, \emph{sigdim}, software, inspected commit:\newline
\path{24f7eb387703fadb0c4d79c831b9def2314fe198}.
\url{https://codeberg.org/mda/sigdim}.

\bibitem{SignalingSoftware}
Chitambar Lab, \emph{SignalingDimension.jl}, software, inspected commit:\newline
\path{5ec10bda0b480e0f811595210d33aad99eba4e0a}.
\url{https://github.com/ChitambarLab/SignalingDimension.jl}.

\bibitem{BenczeFrenkel}
T. Bencze and P. E. Frenkel,
``On the signaling dimension of the cross-polytope,''
arXiv:2610.05924v1 (2026).
\url{https://arxiv.org/abs/2610.05924v1}.

\bibitem{KaiDallArno}
S. Kai and M. Dall'Arno,
``The signaling dimension of two-dimensional and polytopic systems,''
\emph{Quantum Information and Computation} \textbf{24} (2024), 932--944.
\url{https://doi.org/10.26421/QIC24.11-12-3}.

\bibitem{DallArnoTosiniBuscemi}
M. Dall'Arno, A. Tosini and F. Buscemi,
``The signaling dimension in generalized probabilistic theories,''
\emph{Quantum Information and Computation} \textbf{24} (2024), 411--424.
\url{https://doi.org/10.26421/QIC24.5-6-2}.
\url{https://arxiv.org/abs/2311.13103}.

\bibitem{Frenkel}
P. E. Frenkel,
``Classical simulations of communication channels,''
\emph{Quantum} \textbf{6} (2022), 751.
\url{https://doi.org/10.22331/q-2022-06-29-751}.

\bibitem{FrenkelWeiner}
P. E. Frenkel and M. Weiner,
``Classical information storage in an $n$-level quantum system,''
\emph{Communications in Mathematical Physics} \textbf{340} (2015), 563--574.
\url{https://doi.org/10.1007/s00220-015-2463-0}.

\bibitem{NoHypersignaling}
M. Dall'Arno, S. Brandsen, A. Tosini, F. Buscemi and V. Vedral,
``No-hypersignaling principle,''
\emph{Physical Review Letters} \textbf{119} (2017), 020401.
\url{https://doi.org/10.1103/PhysRevLett.119.020401}.

\bibitem{DoolittleChitambar}
B. Doolittle and E. Chitambar,
``Certifying the classical simulation cost of a quantum channel,''
\emph{Physical Review Research} \textbf{3} (2021), 043073.
\newline
\url{https://doi.org/10.1103/PhysRevResearch.3.043073}.

\bibitem{McKayPiperno}
B. D. McKay and A. Piperno,
``Practical graph isomorphism, II,''
\emph{Journal of Symbolic Computation} \textbf{60} (2014), 94--112.
\url{https://doi.org/10.1016/j.jsc.2013.09.003}.

\end{thebibliography}

\end{document}
